\documentclass[../main]{subfiles}
\begin{document}
\chapter{Estudio de Mercado}
\img{images/06/estudio-de-mercado.jpg}{Gráficos de barra y diagramas de torta.}
\info{Estudio de mercado}{La investigación de mercados es la
  herramienta necesaria para la identificación, acopio, análisis,
  difusión y aprovechamiento sistemático y objetivo de la información,
con el fin de mejorar la toma de decisiones relacionadas con la mercadotecnia.}
\section{Clasificación de bienes}
\subsection{Bienes finales}
Son a los que les daremos un uso final, es decir, han sido producidos
directamente para ser utilizados por el consumidor final y satisfacer
una necesidad determinada (por ejemplo, una mesa, un jersey o un lápiz).
\subsection{Bienes Intermedios}
Son todos aquellos recursos materiales, bienes y servicios que se
utilizan como productos intermedios durante el proceso productivo,
tales como materias primas, combustibles, útiles de oficina, etc. Se
compran para la reventa o bien se utilizan como insumos o materias
primas para la producción y venta de otros bienes.
\subsection{Bienes de capital}
Maquinaria y equipo que forman parte del activo fijo de las empresas
y son utilizados para la producción de otros bienes. Es cualquier
bien que se utiliza en un proceso productivo, permitiendo producir
otros bienes, servicios o riquezas.
\section{Mercado Proveedor}
\info{Mercado proveedor}{
  El análisis de proveedores es una herramienta que se utiliza para
  analizar y seleccionar a los proveedores que se necesitan para poner
en marcha un negocio.}
Un proveedor puede ser cualquier persona o empresa que abastece los
recursos que otra empresa necesita para llevar a cabo su proceso de
producción. Los recursos pueden ser bienes y servicios que son
transformados para producir un producto.
\subsection{¿Por qué hacer un análisis de proveedores?}
Encontrar buenos proveedores es vital para una empresa por las
siguientes razones:
\begin{enumerate}
  \item Toda cadena de suministros inicia con un buen proveedor, por
    lo que se vuelve indispensable escoger a los mejores proveedores.
  \item La selección de un buen proveedor hace que cumpla mejor con
    los requerimientos y las necesidades de la empresa.
  \item Una empresa no puede sobrevivir si no cuenta con insumos y,
    como consecuencia, necesita proveedores.
  \item Los proveedores inciden sobre los costos y la calidad de los
    productos que elabora una empresa.
\end{enumerate}
\subsection{Factores importantes en el análisis de proveedores}
Una empresa tiene que realizar un análisis de proveedores por razones
diferentes. Por ejemplo, puede que sea una nueva empresa que inicia
operaciones y no cuenta con proveedores. Otra situación se da cuando
la empresa ya cuenta con proveedores, pero no está satisfecha con los
servicios recibidos, o simplemente quiere ampliar su cartera de proveedores.
Los factores más importantes en el análisis de proveedores son:
\subsubsection{Calidad}
Entre los aspectos relacionados con la calidad encontramos:
\begin{enumerate}
  \item Calidad del producto: Para comprobar la calidad del producto
    se hacen evaluaciones y pruebas sobre las características que son
    consideras importantes. El criterio de calidad se maneja conforme
    a lo que necesita la empresa. Si dos proveedores ofrecen la misma
    calidad, se debe escoger el más económico. Pero, debe quedar
    claro que no siempre lo más barato es lo mejor.
  \item Características técnicas: Las características técnicas del
    proveedor corresponden al equipo y la maquinaria que utiliza para
    elaborar sus productos, de manera que la empresa compradora pueda
    elegir a su mejor proveedor.
  \item Garantía: La empresa siempre buscará encontrar en sus
    proveedores una garantía lo más extendida que sea posible.
  \item Personal técnico: Se analiza que el proveedor cuente con
    personal calificado para asistir a la empresa.
    Asistencia técnica y servicio posventa: Se refiere al tiempo que
    el proveedor otorga para garantizar el mantenimiento, la
    reparación o la asistencia técnica sobre lo que se ha comprado.
\end{enumerate}
\subsubsection{Condiciones económicas}
\img{images/06/proveedores.png}{Análisis de proveedores}
Entre las condiciones económicas se analiza:
\begin{itemize}
  \item Precio: Se busca un precio competitivo que se relacione con
    la calidad ofrecida. Se considera el precio, los descuentos, los
    gastos de transporte, embalajes, carga y descarga.
  \item Forma de pago: La forma de pago analizada puede ser al
    contado, al crédito o conforme a las políticas de la empresa.
  \item Precio de envase y embalaje: El precio del embalaje y envase
    deberá estar acorde con el nivel de protección del producto y que
    sea de costo bajo.
    Pago de seguros: Si el proveedor se hace cargo del pago de seguros.
  \item Recargos por aplazar pagos: Se revisa el recargo que aplica
    el proveedor en el caso de atraso en algún pago.
  \item Descuentos por pronto pago: Descuentos que se otorgan por
    pagos al contado o por períodos menores a los convenidos.
\end{itemize}
\subsubsection{Otras condiciones}
Otras condiciones que también se analizan son:
\begin{enumerate}
  \item Tiempo válido de la oferta.
  \item Condiciones para terminar contratos.
  \item Situación que provoquen revisión de precios.
  \item Plazos de entrega.
  \item Devoluciones.
  \item Inventarios.
  \item Infraestructura.
  \item Legalidad en su funcionamiento.
  \item Tiempo y experiencia de la empresa proveedora.
  \item Recomendaciones de otros compradores.
\end{enumerate}
\section{Mercado del consumidor}
El mercado de consumidores representa todos los compradores que
buscan adquirir los bienes y servicios que se venden en el mercado
para ser usados para satisfacer una necesidad, por ello se les llama
consumidores porque son los que usan y consumen los productos.
\section{Mercado del Proyecto}
\subsection{Competencia}
El mercado competitivo, en el que hay un gran número de compradores y
vendedores, y los vendedores ofrecen productos idénticos a los
compradores; se conoce como competencia perfecta. La competencia
imperfecta ocurre cuando no se cumplen una o más condiciones de la
competencia perfecta.
\subsection{Sustitutos}
Son bienes que satisfacen una misma necesidad o una de muy similar y,
por lo tanto, pueden ser reemplazados por otro bien según los
factores que decanten la decisión del comprador (como el precio).
Ejemplos muy claros de productos sustitutivos son la margarina y la
manteca o el azúcar y el edulcorante.
\section{Mercado global}
\img{images/06/mercadoglobal.png}{Mercado global es el mercado mundial}
El mercado global o mundial es un sistema de relaciones económicas,
mercantiles y financieras, entre estados enlazados por la división
internacional del trabajo. Con el concepto de la división
internacional del trabajo está íntimamente relacionado el concepto de
cooperación internacional, la base de una administración eficiente de
los factores de producción.
\section{Mercados de servicios.}
Están englobados en el sector terciario de la economía de un país
desde una óptica de marketing, nos estamos refiriendo a aquellos
bienes o productos de naturaleza principalmente intangible que
satisfacen la cada vez mayor demanda de este tipo de productos.
\img{images/06/atencion_al_cliente.png}{Servicio de atención al cliente}
\actividad{ Actividad para el proyecto}
\begin{enumerate}
  \item Clasificar el tipo de bien que se trata el producto del
    proyecto. Y justificar la respuesta.
  \item Elaborar una clasificación de los bienes intermedios
    necesarios para producir el producto y clasificar la elasticidad
    del mercado de proveedores en cada caso y también para el
    producto del proyecto.
  \item Elaborar una lista con los sustitutos al producto.
\end{enumerate}
\actividad{ Investigar y contestar}
\begin{enumerate}
  \item ¿Qué es la comercialización de un producto?
  \item ¿Que es el marketing y como se diferencia de la publicidad?
  \item ¿Qué es la mercadotecnia?
  \item ¿Qué es la elasticidad de la demanda de un producto?
  \item ¿Qué significa que un producto tenga demanda elástica?
  \item Buscar ejemplos de productos que tengan un mercado de demanda
    inelástica.
  \item Dar 5 ejemplos de mercados de servicios.
  \item Dar la definición de mercado monopólico, oligopólico y de
    competencia perfecta.
  \item ¿Qué son los factores de la producción?
\end{enumerate}
\end{document}

